\chapter{Fonctions trigonométriques et fonctions réciproques}
\label{workbook:chapter:16}
\section{Graphes fondamentaux}
\label{workbook:section:16.01}
\begin{exerciseblock}{Graphes fondamentaux}
\item Donner domaine, image, période, zéros et parité de $\sin x$, $\cos x$ et $\tan x$.
\item Esquisser les trois graphes sur $[-2\pi,2\pi]$ en marquant les points et asymptotes principaux.
\item Résoudre graphiquement $\sin x=\cos x$ sur $[0,2\pi]$.
\end{exerciseblock}
\section{Amplitude, période, déphasage et ligne médiane}
\label{workbook:section:16.02}
\begin{exerciseblock}{Amplitude, période, déphasage et ligne médiane}
\item Pour $y=3\sin(2(x-\pi/4))-1$, identifier amplitude, période, déphasage et ligne médiane.
\item Construire une équation sinusoïdale de maximum $7$, minimum $1$, période $6$ et maximum à $x=2$.
\item Expliquer l'effet séparé de $A$, $B$, $C$, $D$ dans $A\sin(B(x-C))+D$.
\end{exerciseblock}
\section{Fonctions trigonométriques réciproques}
\label{workbook:section:16.03}
\begin{exerciseblock}{Fonctions trigonométriques réciproques}
\item Calculer exactement $\arcsin(1/2)$, $\arccos(-\sqrt2/2)$ et $\arctan(1)$ dans leurs intervalles principaux.
\item Expliquer pourquoi $\arcsin(\sin(5\pi/6))\neq5\pi/6$.
\item Donner domaine et image de $\arcsin$, $\arccos$ et $\arctan$.
\end{exerciseblock}
\section{Équations trigonométriques élémentaires}
\label{workbook:section:16.04}
\begin{exerciseblock}{Équations trigonométriques élémentaires}
\item Résoudre $\sin x=\sqrt3/2$ sur $[0,2\pi]$ puis dans $\R$.
\item Résoudre $2\cos^2x-1=0$ sur $[0,2\pi[$.
\item Résoudre $\tan(2x)=1$ sur $[0,\pi]$.
\end{exerciseblock}
\section{La sinusoïde est l'ombre d'une rotation}
\label{workbook:section:16.05}
\begin{exerciseblock}{La sinusoïde est l'ombre d'une rotation}
\item Paramétrer un point tournant sur un cercle de rayon $4$ et montrer que son ordonnée est sinusoïdale.
\item Associer les points clés du cercle aux points $0,\pi/2,\pi,3\pi/2,2\pi$ du graphe du sinus.
\item Expliquer pourquoi la projection horizontale donne le cosinus et la verticale le sinus.
\end{exerciseblock}
\section{Effet des paramètres d'une sinusoïde}
\label{workbook:section:16.06}
\begin{exerciseblock}{Effet des paramètres d'une sinusoïde}
\item Comparer sur un même repère $\sin x$, $2\sin x$, $\sin(2x)$ et $\sin(x-\pi/2)$.
\item Pour $f(x)=-4\cos(3x)+2$, déterminer maximum, minimum, période et premières positions de maximum.
\item Trouver les paramètres d'une sinusoïde à partir d'un graphe ayant maximum $5$, minimum $-1$, période $4\pi$ et passage croissant par la médiane à $x=\pi$.
\end{exerciseblock}
\section{Construire une sinusoïde et retrouver un angle}
\label{workbook:section:16.07}
\begin{exerciseblock}{Construire une sinusoïde et retrouver un angle}
\item À partir de maximum, minimum, période et position initiale, donner une procédure de construction en cinq points.
\item Construire $y=2\sin(\pi t/3)-1$ sur une période et marquer les quarts de période.
\item Résoudre $2\sin(\pi t/3)-1=0$ sur une période, puis interpréter les deux instants.
\end{exerciseblock}
\section*{Synthèse du chapitre}
\begin{synthesisbox}
\item Étudier complètement $f(x)=2\cos(3(x-\pi/6))+1$ : paramètres, points clés, zéros et solutions de $f(x)=2$ sur une période.
\item Un graphe sinusoïdal a maximum $12$ à $t=1$, minimum $4$ à $t=5$, puis maximum suivant à $t=9$. Construire un modèle cosinus.
\item Résoudre $\sin x+\cos x=1$ sur $[0,2\pi]$ en utilisant une identité ou une transformation de phase.
\end{synthesisbox}
